\section{Conclusion}
\label{sec:conclusion}

We introduced a publication-derived benchmark of 40 PaN papers, with reconciled 
gold records and a two-layer schema for generic and experimental metadata. 
Its entity-aware evaluation aligns entities and remaps output-local identifiers, 
separating entity recovery from conditional attribute and reference quality.

Both annotators and models generally recover paper-level fields, authors,
and affiliations more consistently than facilities, samples, and datasets. 
PaN annotation disagreements frequently concern entity identity and granularity. 
No model leads all three PaN collections, and higher entity recovery does not 
necessarily imply higher conditional attribute or reference quality. Detailed field 
instructions improve dataset and sample recovery across all evaluated models, but 
their effects on facilities are model-dependent. Within the Qwen3.5 family, the 
association between model size and entity recovery depends on the entity
type, while fixed demonstrations produce heterogeneous and sometimes
non-monotonic changes. Together, these findings motivate reporting entity
recovery, attribute quality, and reference quality separately instead of
collapsing them into a benchmark-wide score.

We release the schema, gold records, scoring protocol, and evaluation code to 
support further study of these distinct extraction challenges.
